\documentclass[11pt]{article}
\usepackage[a4paper,margin=25mm]{geometry}
\usepackage{amsmath,amssymb,amsthm,mathtools,booktabs,hyperref,xurl}
\setlength{\emergencystretch}{2em}
\newtheorem{theorem}{Theorem}
\newtheorem{lemma}[theorem]{Lemma}
\newtheorem{proposition}[theorem]{Proposition}
\newtheorem{corollary}[theorem]{Corollary}
\newcommand{\Z}{\mathbb Z}\newcommand{\Q}{\mathbb Q}
\title{A single integer quotient for the third Pell certificate}
\author{\(\mathbb Z\mathbb Q\) Diophantine Project}
\date{23 November 2026}
\begin{document}
\maketitle

\begin{abstract}
The third Pell equation in the corrected Jones source used two integer
quotients: one fixed its second coordinate modulo the first Pell
coordinate, and another fixed its trace modulo an auxiliary Pell
trace. We replace these by two congruences on the second coordinate
alone, and encode both congruences with one integer quotient. An
elementary signed congruence lemma for the Pell second-coordinate
sequence proves soundness. Every positive source witness has a new
witness with the quotient positive. One new quadratic product register
replaces the discarded source coordinate, so the outer system has the
same size and one entire integer predicate can be removed. The general,
specified, and ordinary rational-\(abc\) complete tables decrease to
983, 1148, and 734 coordinates, respectively, with 855, 998, and 643
quadratic residuals. This local algebraic and Pell
reduction is unconditional; the undecidability transfer retains the
stated arithmetic premises of its integer interface and the cited
Jones coding theorem.
\end{abstract}

\section{A second-coordinate congruence lemma}

For an integer \(A\ge2\), write
\[
 X_n(A)+Y_n(A)\sqrt{A^2-1}
       =(A+\sqrt{A^2-1})^n,\qquad n\ge0.
\]
Thus \(X_0=1,Y_0=0,X_1=A,Y_1=1\), and both sequences
satisfy \(Z_{n+2}=2AZ_{n+1}-Z_n\). For negative integer
parameters we use these same integer polynomials; then
\[
 X_n(-A)=(-1)^nX_n(A),\qquad
 Y_n(-A)=(-1)^{n-1}Y_n(A).
\]
All congruences below use a positive modulus unless explicitly stated.

\begin{lemma}[The signed \(Y\)-congruence]\label{Ylemma}
Let \(A\ge2\), \(q\ge2\), and \(0<t<q\) be integers. For
every integer \(v\ge0\),
\[
 Y_v(A)\equiv\pm Y_t(A)\pmod{X_q(A)}
 \quad\Longrightarrow\quad
 v\equiv\pm t\pmod{2q}.
\]
\end{lemma}
\begin{proof}
Suppress the argument \(A\) and put \(N=X_q\). The Pell addition
identities and the norm identity give
\[
 X_{2q}=2X_q^2-1\equiv-1\pmod N,
 \qquad Y_{2q}=2X_qY_q\equiv0\pmod N.
\]
Consequently
\[
 Y_{v+2q}\equiv-Y_v\pmod N,
 \qquad Y_{2q-j}\equiv Y_j\pmod N\quad(0\le j\le2q).
\]
Reduction modulo \(4q\), followed by reflection around \(q\),
therefore gives some \(j\in\{0,\ldots,q\}\) such that
\[
 Y_v\equiv\varepsilon Y_j\pmod N,
 \qquad v\equiv\pm j\pmod{2q},\qquad\varepsilon\in\{1,-1\}.
\]

The sequence \(Y_n\) is strictly increasing for \(n\ge0\), and
\(Y_q>2Y_{q-1}\) for \(q\ge2\). For example the recurrence
gives \(Y_2=2A>2Y_1\), and if \(Y_j>Y_{j-1}\) then
\(Y_{j+1}=(2A-1)Y_j+(Y_j-Y_{j-1})>2Y_j\).
The identity \(X_q=AY_q-Y_{q-1}\) now yields
\[
 X_q-(Y_q+Y_{q-1})=(A-1)Y_q-2Y_{q-1}>0.
\]
Since \(t\le q-1\), we have
\[
 0<Y_j+Y_t\le Y_q+Y_{q-1}<N.
\]
Thus \(Y_j\equiv-Y_t\pmod N\) is impossible, and
\(Y_j\equiv Y_t\pmod N\) implies equality because both values
lie in \([0,N)\). Strict monotonicity gives \(j=t\). Substituting
this in the congruence for \(v\) proves the lemma.
\end{proof}

\section{The third norm with one integer quotient}

\begin{theorem}[A one-quotient index certificate]\label{index}
Let \(a,c,d,B,i,J\) be integers satisfying
\[
 |a|\ge2,\qquad B\ge3\text{ odd},\qquad |c|>B,\qquad i\ne0.
\]
Let \(F,I\in\Q\), and suppose
\begin{align}
 d^2-(a^2-1)c^2&=1,\label{norm1}\\
 F^2-(a^2-1)i^2c^4&=1.\label{norm2}
\end{align}
Set
\begin{align}
 G&=a+F^2(c+1-a),\label{bridge}\\
 H&=c+F^2(B-c)+cFJ.\label{CRT}
\end{align}
Assume
\begin{equation}\label{norm3}
                       I^2-(G^2-1)H^2=1.
\end{equation}
Then
\[
              |c|=Y_B(|a|),\qquad |d|=X_B(|a|).
\]
If a rational \(\zeta\) also satisfies
\(\zeta^2=(a-1)(d-1)\), then \(d=X_B(a)\).
\end{theorem}
\begin{proof}
Put \(A=|a|\) and \(C=|c|\). Integral Pell classification
in \eqref{norm1} gives
\[
 C=Y_t(A),\qquad |d|=X_t(A)
\]
for an integer \(t\ge2\), since \(C>B\ge3\). In particular
\(C\ge2A\) and \(0<t<C\). A rational number whose square is
an integer is an integer, so \eqref{norm2} makes \(F\) integral.
Pell classification gives
\[
 |i|C^2=Y_q(A),\qquad |F|=X_q(A)
\]
for an integer \(q\ge1\). The square-divisibility step-down
lemma proved in \cite{Pell} states that
\(Y_t(A)^2\mid Y_q(A)\) implies \(Y_t(A)\mid q\).
Therefore \(C\mid q\), whence \(q\ge C>t\), in particular
\(q\ge2\). Also
\begin{equation}\label{Fcong}
 F^2\equiv1\pmod C,\qquad \gcd(C,F)=1.
\end{equation}

All coordinates in \eqref{bridge}--\eqref{CRT} are now integers.
The bridge is nondegenerate for every sign of \(a,c\). If
\(c=C>0\), then \(c+1-a\ge C+1-A>0\) and \(F^2\ge1\)
give \(G\ge c+1\ge5\). If \(c=-C<0\), then
\(c+1-a\le1-A<0\), so \(G\le c+1\le-3\).
Thus \(|G|\ge2\). Moreover
\[
 G\equiv1\pmod C,\qquad G\equiv a\pmod {|F|}.
\]
Equations \eqref{CRT} and \eqref{Fcong} give the two congruences
\begin{equation}\label{bothcong}
                    H\equiv B\pmod C,
                    \qquad H\equiv c\pmod {|F|}.
\end{equation}
Since \(0<B<C\), the first congruence implies \(H\ne0\).
The last norm makes \(I\) an integer. Signed Pell classification
at \(G\) consequently gives
\[
                  H=\varepsilon Y_v(G),\qquad v\ge1,
                  \quad\varepsilon\in\{1,-1\}.
\]

The recurrence at parameter \(1\) gives \(Y_v(1)=v\).
The first congruence in \eqref{bothcong} therefore yields
\(v\equiv\pm B\pmod C\). Reducing the second congruence
modulo \(|F|=X_q(A)\), using \(G\equiv a\) and the sign
identity for \(Y_v(-A)\), gives
\[
                   Y_v(A)\equiv\pm Y_t(A)\pmod{X_q(A)}.
\]
Lemma~\ref{Ylemma} implies \(v\equiv\pm t\pmod{2q}\).
Because \(C\mid q\), we obtain \(B\equiv\pm t\pmod C\).
Both \(B\) and \(t\) lie strictly between \(0\) and \(C\),
so either \(B=t\) or \(B=C-t\). The recurrence modulo two
gives \(C=Y_t(A)\equiv t\pmod2\); hence \(C-t\) is even,
while \(B\) is odd. The second case is impossible and \(t=B\).

Finally \(|d|=X_B(A)>1\). If \(a>0\), the square equation
for \(\zeta\) forces \(d\ge1\); if \(a<0\), it forces
\(d\le1\). Since \(B\) is odd, this says precisely
\(d=X_B(a)\).
\end{proof}

\begin{proposition}[Positive completeness]\label{complete}
Suppose \(a\ge2\), \(B\ge3\) is odd,
\(c=Y_B(a)>B\), \(d=X_B(a)\), and positive integers \(i,F\)
satisfy \eqref{norm2}. Then positive integers \(J,H,I\) satisfying
\eqref{bridge}--\eqref{norm3} exist.
\end{proposition}
\begin{proof}
As in Theorem~\ref{index}, write \(F=X_q(a)\) with \(c\mid q\).
Then \(G\ge c+1\ge2\), \(G\equiv1\pmod c\), and
\(G\equiv a\pmod F\). For any sufficiently large integer
\(L\ge0\), put
\[
 v=B+4qL,\qquad H=Y_v(G),\qquad I=X_v(G).
\]
The Pell norm is automatic. Reduction modulo \(c\) gives
\(H\equiv v\equiv B\pmod c\). The identities used in
Lemma~\ref{Ylemma} give period \(4q\) for \(Y_v(a)\) modulo
\(F\), hence
\[
 H\equiv Y_v(a)\equiv Y_B(a)=c\pmod F.
\]
Because \(F^2\equiv1\pmod c\), the integer
\(H-c-F^2(B-c)\) is divisible by both \(c\) and \(F\).
They are coprime, so
\[
                   J=\frac{H-c-F^2(B-c)}{cF}
\]
is an integer. As \(L\) grows, \(H\) tends to positive infinity
while the other terms are fixed. Thus \(J>0\) for all sufficiently
large \(L\), establishing the claim including its positivity.
\end{proof}

\section{Compatibility with the signed binomial block}

In the corrected source of \cite{Corrected}, and after the signed
positivity reduction of \cite{Signed}, the relevant parameters are
\[
 M=n^2\ge16,\quad r\ge1,\quad B=2r+1,\quad
 Y=sM,\quad a=(wM+1)rY,\quad p=2wr^2Y^2M,
\]
where \(s\ne0\), \(w,s\in\Z\), and \(2\le b\le2M\).
The small Pell equation, its integer congruence quotient, and
\(k^2-4(c-kY)^2>0\) are unchanged. The proof of the automatic
positivity theorem first excludes \(w=0\), then establishes
\(|a|\ge2\), \(|c|>B\), and the integrality of every input needed
by Theorem~\ref{index}. Those arguments precede, and do not use,
the old third-Pell index lemma. Substituting Theorem~\ref{index}
for that lemma gives exactly the same conclusion
\[
                 |c|=Y_B(|a|),\qquad d=X_B(a).
\]
The rest of the signed proof uses these conclusions, the unchanged
linear power-of-two congruence, and the small Pell equation; it does
not use either of the discarded third-norm quotient equations.
Therefore it again forces \(w,s,c>0\),
\(bw=2^{2r+1}\), and \(M\mid\binom{2r}{r}\).

The reciprocal guard has the same noncircular order as before.
Since \(s\ne0\) and \(M\ge16\), the strict inequality implies
\(k,c\ne0\). The retained equation \(vE=c^2\) forces \(v\ne0\),
so the integer predicate for \(1/v\) gives
\(i\in\Z\setminus\{0\}\); the second norm then gives integral
\(F\). In the ordinary-\(abc\) interface the explicit row
\(vi=1\) is retained as well. No new nonzero assumption is implicit.

Conversely, the corrected positive source has its positive local
Pell witnesses by the established binomial certificate.
Proposition~\ref{complete} supplies a positive new quotient \(J\).
Thus all eleven remaining integer inputs can be taken nonzero, as
required by the fixed-contact ordinary-\(abc\) completeness theorem.
The new system has the same projection on the coding-prefix inputs
and on the instance \(x\). This is existential equivalence of the
source projections: the integer quotient named \(j\) in the new table
means \(J\), and it need not equal the old quotient named \(j\).

\section{Complete quadratic compilation and counts}

The complete predecessor tables already store
\(F_2=F^2\), \(r,c,F,H,I,G\). Introduce the single new register
\(C_F=cF\), delete the source coordinate \(o\), and replace
\[
                       H=2r+1+jc,\qquad I=d+oF
\]
by the two quadratic rows
\begin{equation}\label{tworows}
            C_F=cF,\qquad H=c+F_2(2r+1-c)+C_Fj.
\end{equation}
Retain the third norm and all other source rows, and interpret the
integral input \(j\) as the new quotient \(J\). There is one fewer
source coordinate and one new product coordinate, so the outer
coordinate count is unchanged. Two old residuals are replaced by two
new residuals, so the outer row count is also unchanged.

The ordinary-\(abc\) predecessor's \(o\)-input certificate has
47 private auxiliary coordinates and 41 private quadratic residuals.
After the row \(I=d+oF\) is removed, none of those coordinates,
nor \(o\) itself, occurs in any retained source or other certificate
row. Deleting this certificate and adding \(C_F\) therefore changes
the total variable count by \(-48+1=-47\), and the row count by
\(-41\). All eleven retained fixed contact offsets are still distinct
modulo \(\Z_2\); the removed offset was \(9/16\), and keeping the
other original offsets causes no collision. The common reference
certificate and common contact square remain. The simultaneous
completeness theorem for a finite family applies directly to the
eleven retained inputs; no unused twelfth integer witness is needed.

\begin{corollary}[A 734-coordinate ordinary-\(abc\) quartic]\label{abc}
Under exactly the arithmetic hypotheses of the retained ordinary
rational-\(abc\) integer interface, and using the cited Jones coding
theorem as in \cite{Corrected}, the new complete family has
undecidable rational solvability. Its instance is the coefficient
\(x\), its fixed source parameters are \(z,u,y\), and it has
734 rational unknowns and 643 quadratic residuals.
Squaring and summing those residuals gives a polynomial of degree
exactly four, with coefficients in \(\Z[x,z,u,y]\).
\end{corollary}
\begin{proof}
Soundness of the eleven retained integer predicates and
Theorem~\ref{index}, in the signed reconstruction just described,
recover a solution of the corrected coding source. Conversely a
positive corrected source witness extends by
Proposition~\ref{complete}, and finite-family simultaneous
completeness supplies the eleven rational integer certificates.
The coefficient table makes the two row changes in \eqref{tworows},
deletes exactly the private certificate just counted, and preserves
all denominator guards. Thus the counts are
\(781-47=734\) and \(684-41=643\). Every retained or added row
has degree at most two, and the coefficient domain is unchanged.
A sum of rational squares vanishes if and only if all its terms
vanish. At least one row has a nonzero quadratic homogeneous part;
the sum of the squares of these parts cannot vanish identically over
the reals. Thus the degree is exactly four.
\end{proof}

\section{Deleting one member of the simultaneous integer family}
\label{sec:delete-integer-family-member}

The new CRT bridge needs the integer quotient $j$, now interpreted as
$J$, but no longer needs the old integer $o$. The general and specified
rational interfaces can therefore be instantiated with eleven input
blocks. This is a new application of their finite-family theorem; it
does not assert that a rational zero of an eleven-block system extends
to a zero of the former twelve-block system with a prescribed $o$.

Here is the required family argument. Keep the blocks indexed by
\[
              0,1,2,3,4,5,6,7,8,10,11,
\]
with the original ordered labels $b,g,t,m,w,s,h,c,\gamma,j,i$.
The reference block $b$, common point roles, common input numerator,
common contact gauge, and reciprocal $i$ chart remain. For $e=2$
general contacts or $e=3$ specified contacts, define the gauge by the
sum of the retained weights only. Thus its number of summands is now
$m=22$ or $m=33$, respectively.

The soundness proof of the global gauge uses a consumed prime of one
retained weight, its own complementary sum, the unchanged scalar and
point relations, and its own pole obstruction. It makes no assumption
that another input block exists. All remaining height estimates use
the same finite-family inequalities, with fewer summands; the former
twelve-block upper bounds are still sufficient. Hence every remaining
distinguished input is an integer under the same arithmetic premises.
The affine contact chart is reconstructed with the seven common
coefficient registers exactly as before.

For completeness, start with the eleven desired integer inputs and
choose their point data and aggregate shifts by the existing ordered
simultaneous preparation. Take the common contact spacing sufficiently
large to separate the translated finite lists. The displayed indices
are distinct integers, so the interval-separation proof applies even
though index $9$ has been omitted. Invoke the same finite-family prime
supplier on their union, with the same finite avoidance and all-sign
common-index congruences. At each selected contact, only its own weight
has negative odd valuation, and all other retained weights are integral.
At foreign contacts, the same rational-square parity argument applies.
At the prepared finite places and at infinity,
\[
 q_\nu=\kappa\lambda^{-2}+O(\lambda^{-3}),\qquad
 \rho=m\kappa\lambda^{-2}+O(\lambda^{-3}),\qquad
 \frac{(\rho-q_\nu)^2}{q_\nu}
       =(m-1)^2\kappa\lambda^{-2}+O(\lambda^{-3})\longrightarrow0.
\]
No division by $m$ or $m-1$ is needed. The physical tests are bounded
on the previously chosen neighborhoods; their nonconstant detector
terms tend to zero. The actual pole arguments thus retain the required
local square classes at $2,3,17$ and the same real limit. Choosing
sufficient finite divisibility and a large real slope gives every
remaining native-pole witness. This proves the eleven-block
simultaneous completeness assertion.

These are precisely the arguments in
\emph{One contact gauge in a combined quadratic system},
\nolinkurl{undecidability/papers/global_gauge_quartic_checkpoint.tex},
Sections 2 and 3, and
\emph{A common arithmetic progression of contact arguments},
\nolinkurl{undecidability/general_contact_v94/affine_contact_progressions.tex},
with a smaller finite family. No new prime-supplier premise, positivity
assumption on an arbitrary rational zero, or pairwise radical
disjointness assumption is introduced.

\paragraph{Complete coefficient audit.}
In the signed affine-contact tables, the literal private coordinates
are $o$ and every name beginning \texttt{int\_o\_}. There are $61$
of them in the general table and $76$ in the specified table. The
private integer-certificate rows number $46$ and $59$. Each interface
also has four private \texttt{global.o.} input-pair rows, one private
\texttt{scalar\_norm.o.} row, and one private
\texttt{detector\_trace.o.} row. Their deletion removes $52$ and
$65$ quadratic rows, respectively.

Apart from the old source row $I-d-oF$, the only occurrences of these
private coordinates outside those rows lie in three rows of the last
reciprocal-$i$ contact. They are exactly the eliminated global-sum
contributions of the removed weights. The general contributions are
\texttt{int\_o\_153} and
\texttt{int\_o\_last\_contact\_weight}; the specified contributions
are \texttt{int\_o\_154}, \texttt{int\_o\_155}, and
\texttt{int\_o\_last\_contact\_weight}.
Deleting their terms from these three rows compiles the smaller global
sum. All other retained rows, including the common gauge guards and
the private guards of every retained block, are literally unchanged.

The CRT source then drops $I-d-oF$, replaces the old definition of $H$,
and adds one coordinate $CF$ with the equation $CF-cF=0$. Its new
quadratic $H$ equation is
\[
              H-c-2rF_2-F_2+cF_2-(CF)j=0,
              \qquad F_2=F^2.
\]
Thus the source row count is unchanged, and its additional coordinate
offsets one of the deleted private coordinates. The full counts become
$1043-61+1=983$ variables and $907-52=855$ rows in the general
interface, and $1223-76+1=1148$ variables and $1063-65=998$ rows in
the specified interface. The exact coefficient checker independently
verifies every private support, every deleted global-sum term, every
unchanged row, the new source equations, coefficient domains, guards,
degree, counts, and complete coordinate support.


\section{Exact artifacts and scope of verification}

The independent scripts
\begin{verbatim}
python undecidability/pell_crt_v95/verify_y_congruence.py
python undecidability/pell_crt_v95/verify_crt_abc.py
python undecidability/pell_crt_v95/verify_crt_general.py
\end{verbatim}
check the recurrence identities, signed congruence examples, positive
CRT lifts, every coefficient occurrence needed for deletion, every
unchanged row, counts, degree, coefficient hashes, and deterministic
rebuilding. Finite examples do not prove the quantified Pell theorem;
its proof is the argument above. These checks neither reprove the
external Jones coding theorem nor discharge any arithmetic premise
of the integer interfaces. No global minimum or 263-variable quartic
is claimed.

\begin{thebibliography}{9}
\small
\setlength{\itemsep}{2pt}
\bibitem{Pell} \(\mathbb Z\mathbb Q\) Diophantine Project,
\emph{An exact three-Pell index certificate and nine fewer quadratic
coordinates}, 23 November 2026.\par
\nolinkurl{undecidability/papers/pell_index_optimization.tex}.
\bibitem{Corrected} \(\mathbb Z\mathbb Q\) Diophantine Project,
\emph{A corrected quartic source from Jones's binomial theorem},
23 November 2026.\par
\nolinkurl{undecidability/source_repair_v93/corrected_jones_source.tex}.
\bibitem{Signed} \(\mathbb Z\mathbb Q\) Diophantine Project,
\emph{Signed binomial witnesses force their own positivity and save
five quadratic coordinates}, 23 November 2026.\par
\nolinkurl{undecidability/source_repair_v93/signed_binomial_positivity.tex}.
\end{thebibliography}
\end{document}
